\subsection{双对偶}

定义 1. --- 设 E 是一个局部凸空间，\(E_{b}^{\prime}\) 是它的强对偶。局部凸空间 \(E_{b}^{\prime}\) 的对偶称为 E 的双对偶，记作 \(E^{\prime\prime}\)。

对每个 \(x \in E\)，设 \(\tilde{x}\) 是线性型 \(x' \mapsto \langle x, x' \rangle\)（\(E'\) 上的）；它对于弱拓扑 \(\sigma(E', E)\) 连续，故更加对于 \(E'\) 上的强拓扑连续；因此 \(\tilde{x} \in E''\)（对所有 \(x \in E\)）。映射 \(c_{E}: x \mapsto \tilde{x}\)（从 E 到 \(E''\) 中的）是一个线性映射，称为典范映射。

命题 1. --- \(c_{\mathrm{E}}: \mathrm{E} \to \mathrm{E}''\) 的核是 E 中 0 的闭包。若 E 是 Hausdorff 的，则 \(c_{\mathrm{E}}\) 是单射。

由构造，\(c_{\mathrm{E}}\) 的核是 E 上连续线性型的核的交，即 \(\{0\}\) 在 E 中的闭包（II，第 24 页，推论 1）。

当 E 是 Hausdorff 的时候，我们把 E 等同于 \(E''\) 的一个子空间，借助的是映射 \(c_{E}\)。

\(E''\) 上的强拓扑是 S-拓扑，其中 S 是 \(E'\) 中所有强有界子集组成的族。由于 \(E'\) 的每个等度连续子集都强有界（III，第 22 页，命题 9），E 上的初始拓扑比在 \(c_{E}\) 下取 \(E''\) 上强拓扑的逆像所得到的拓扑粗；它可能严格地粗（IV，第 52 页，习题 1）。然而：

命题 2. --- 假定空间 E 是有界型的或桶型的。E 上的初始拓扑是 \(c_{E}\) 下 \(E''\) 上强拓扑的逆像。

因为 E' 中每个强有界的子集都等度连续（III，第 22 页，命题 10 与 III，第 24 页）。

命题 3. --- 设 E 是一个局部凸 Hausdorff 空间。为使 E 的强对偶 \(E_{b}^{\prime}\) 是桶型的，必要且充分的是：\(E^{\prime\prime}\) 中每个对于 \(\sigma(E^{\prime\prime}, E^{\prime})\) 有界的子集都包含在 E 的某个有界子集对于 \(\sigma(E^{\prime\prime}, E^{\prime})\) 的闭包中。

\(E''\) 的等度连续子集就是包含在（对于 \(E''\) 与 \(E'\) 之间的对偶而言）\(E''\) 的子空间 E 的某个有界子集的双极集中的那些子集。于是只需应用双极定理（II，第 45 页，推论 3）以及桶型空间的定义（III，第 24 页）。

注. --- 设 E 是一个局部凸 Hausdorff 空间，\(E'\) 是它的对偶，\(E''\) 是它的双对偶。我们有 \(E \subset E'' \subset E'^{*}\)，其中 \(E'^{*}\) 是 \(E'\) 的代数对偶。若 B 是 E 的一个有界子集，则其闭包 \(\overline{B}\)（在 \(E'^{*}\) 中，赋以 \(\sigma(E'^{*}, E')\)）包含在 \(E''\) 中：因为 B 的极集 \(U = B^{\circ}\)（在 \(E'\) 中的）是 \(E_b'\) 中 0 的一个邻域，且我们有 \(\overline{B} \subset U^{\circ} \subset E''\)。

